\documentclass{amsart}
\usepackage{amsmath}
\usepackage{amssymb}
\usepackage[mathscr]{eucal}

% --- Other utilities ---
\usepackage[shortlabels]{enumitem}
\usepackage{graphicx}
\usepackage[all]{xy}
\usepackage[dvipsnames]{xcolor}
\usepackage{tikz-cd}

% --- Hyperref should be loaded last ---
\usepackage[colorlinks=true,citecolor=red,urlcolor=blue,linkcolor=blue]{hyperref}

% --- Theorem environments ---
\newtheorem{theorem}{Theorem}[section]
\newtheorem{proposition}[theorem]{Proposition}
\newtheorem{lemma}[theorem]{Lemma}
\newtheorem{remark}[theorem]{Remark}
\newtheorem{definition}[theorem]{Definition}
\newtheorem{example}[theorem]{Example}
\newtheorem{corollary}[theorem]{Corollary}

% --- Commands ---
\newcommand{\thmref}[1]{Theorem~\ref{#1}}
\newcommand{\thmmref}[1]{Theorem$^\prime$~\ref{#1}}
\newcommand{\secref}[1]{Section~\ref{#1}}
\newcommand{\lemref}[1]{Lemma~\ref{#1}}
\newcommand{\propref}[1]{Proposition~\ref{#1}}
\newcommand{\propnref}[1]{Proposition$'$~\ref{#1}}
\newcommand{\corref}[1]{Corollary~\ref{#1}}
\newcommand{\remref}[1]{Remark~\ref{#1}}
\newcommand{\defref}[1]{Definition~\ref{#1}}
\newcommand{\subsecref}[1]{Subsection~\ref{#1}}
\newcommand{\homeo}{\mathrm{Homeo}}
\newcommand{\gal}{\mathrm{Gal}}
\newcommand{\bz}{\mathbb{Q}}
\newcommand{\bn}{\mathbb{N}}
\newcommand{\bq}{\mathbb{Q}}
\newcommand{\br}{\mathbb{R}}
\newcommand{\bc}{\mathbb{C}}
\newcommand{\bp}{\mathbb{P}}
\newcommand{\bs}{\mathbb{S}}
\newcommand{\co}{\mathcal{O}}
\newcommand{\rank}{\mathrm{Rank}}
\newcommand{\cl}{\mathcal{L}}
\newcommand{\lr}{\longrightarrow}
\renewcommand{\hom}{\mathrm{Hom}}
\newcommand{\wt}{\widetilde}
\newcommand{\im}{\mathrm{Im}}
\newcommand{\re}{\mathrm{Re}}
\newcommand{\Span}{\mathrm{span}}
\newcommand{\tor}{\mathrm{Tor}}
\newcommand{\pic}{\mathrm{Pic}}
\newcommand{\flag}{\mathrm{Flag}}
\newcommand{\ul}{\underline}
\newcommand{\fix}{\mathrm{Fix}}

\title[Graded Endomorphisms of $H^*(\mathbb{S}^m \times \mathbb{C}G_{n,k};\mathbb Q)$]{Rational Cohomology Endomorphisms of product of Sphere with Grassmannian and coincidence theory}

\author[M. Mandal]{Manas Mandal}
\address{Indian Institute of Technology Kanpur, Kanpur 208016, India}
	\email{manasm.imsc@gmail.com}
	
	\author[D. Setia]{Divya Setia}
	\address{Institute of Mathematics, Polish Academy of Sciences, Krak{\'o}w 31-027, Poland}
	\email{divyasetia01@gmail.com}
	
	\subjclass[2020]{Primary: 55S37, 08A35,  55M20; Secondary: 14M15}
	\keywords{cohomology endomorphisms, complex Grassmann manifolds, generalized Dold spaces,  fixed point theory, coincidence theory}

\begin{document}
\begin{abstract}
We classified graded endomorphisms of the rational cohomology algebra of the product of a sphere and a complex Grassmannian, whose images are nonzero  in the second cohomology of the Grassmannian.

We also derive necessary conditions for the generalized Dold spaces to satisfy the coincidence property, in particular the fixed-point property. As an application of our results, we obtain several sufficient conditions for the existence of a point of coincidence between a pair of continuous functions on certain generalized Dold spaces. 
\end{abstract}
	
\maketitle

	\begin{proof}
		Let us consider the sum $\sum_{i=0}^d d_{2i}\lambda^i$, when $\lambda$ is an integer. 
		Clearly, $$\sum_{i=0}^d d_{2i}\lambda^i \equiv 1 \pmod{\lambda}.$$ %(see Theorem 2 in \cite{glover-homer}). 
		Hence, $\sum_{i=0}^d d_{2i}\lambda^i \neq 0$, if $\lambda \neq \pm 1$. When $\lambda = 1$, the sum is also positive and therefore nonzero. It remains to consider the case where $\lambda = -1$.
		Let $\chi(\mathbb{R}G_{n,k})$ denote the Euler-Poincar{\'e} characteristic of $\mathbb{R}G_{n,k}$ and be defined by $$\chi(X):=\sum_{i\ge0}\dim H^i(\mathbb{R}G_{n,k};\mathbb{Z}_2)$$ where $\mathbb{R}G_{n,k}$ denotes the Grassmannian of real $k$-planes in $\mathbb{R}^n$.
		Now we observe that 
		$\sum_{i=0}^d d_{2i}(-1)^i = \chi(\mathbb{R}G_{n,k})$ where $d_{2i} = \operatorname{dim} H^{2i}(\mathbb{C}G_{n,k}; \mathbb{Q}) = \dim H^{i}(\mathbb{R}G_{n,k}; \mathbb{Z}_2)$. 
		It is a well known fact that $\chi(\mathbb{R}G_{n,k}) \neq 0$ if $k(n-k)$ is even. 
		
		Let us move to the other case where $\lambda\in \mathbb{Q}\backslash \mathbb{Z}$. Suppose $\sum_{i=0}^{d} d_{2i}\lambda^i =0$ for some $\lambda = \frac{p}{q}$ where $p$ and $q$ are coprime integers. Since $d_0 =d_d =1$, using the rational root theorem $p|1$ and $q|1$. Hence, $\lambda = \pm 1$, which is a contradiction. Therefore, we conclude that $\sum_{i=0}^{d} d_{2i}\lambda^i \neq 0$ for all $\lambda\in \mathbb Q$. 
		%Thus, the sum  $\sum _{i=1}^d d_{2i}\lambda^i$ is nonzero for every integer value of $\lambda.$ 
		\iffalse	Now, suppose for contradiction that there exists a rational number $\lambda = a/b$  written in lowest terms (i.e., $\gcd(a,b)=1$) such that the sum vanishes. Note that in the polynomial $\sum_{i=0}^{d} d_{2i}\lambda^i$, both the leading coefficient $d_d$ and the constant term $d_0$ are equal to $1$. By the Rational Root Theorem, the rational root $\lambda = a/b$ must satisfy $a \mid 1$ and $b \mid 1$, hence $\lambda = \pm 1\in \mathbb Z$. Since this is not possible, \fi
	\end{proof}
\end{document}
